% SPDX-License-Identifier: Apache-2.0
% covers: Ddr3Per
\datasheet{Ddr3Per}{The board's DDR3 memory as an AXI peripheral: the Wishbone bridge and AMD's MIG controller, which makes the design's clock.}

\dsfacts{\code{ddr3/src/lib.rs}}{\code{ddr3::Ddr3Per}}%
{\code{//docs:vreteno}, \code{//docs:axi}}

\subsection*{Function}

\code{Ddr3Per} puts the DDR3 memory of the Alinx AX7A200 on an AXI
link. It is a unit of units with two children. \code{bridge} is
\code{AxiWb} of \code{txhdl\_parts}, which serves a burst a word at a
time as requests on a pipelined Wishbone bus. \code{ctl} is
\code{Ddr3}, AMD's MIG 7 Series controller, which the build generates
as \code{//ddr3:ddr3\_mig}, behind the wrapper
\code{ddr3/hdl/ddr3\_wb32.v}, which gives it a Wishbone of 32-bit
words. Nine Wishbone signals join the two.

\code{Ddr3} is a foreign unit. Its netlist is an instance of the
module \code{ddr3\_wb32}, which the lowering names and does not write.
Its \code{run} is a model for simulation: a Wishbone memory that
stalls while it calibrates. The model does not drive the memory's
pins.

On the link side the peripheral is the client of an \code{AxiPer},
and the link is one port, \code{bus}, a \code{PerPort}:
\code{bus\_req} and \code{bus\_w} in, \code{bus\_ans} and
\code{bus\_r} out. On the
memory side its ports are the board's clock and the controller's
reset in, the design's clock and its reset out, since the controller
makes them, the \code{calib} output, and the memory's pins, with
\code{dq}, \code{dqs} and \code{dqs\_n} as pads both ways.

\subsection*{Parameters}

None. The controller calibrates in hardware, and its simulation
library carries the variant with the fast calibration, so neither
wants a parameter. The bridge inside is
\code{AxiWb<32, 2, 28>}, with \code{AW} equal to 28.

\subsection*{Address map}

The bridge's Wishbone address is a word address of \code{AW}, 28
bits: the burst's byte address shifted down by two and cut to 28 bits.
\code{lib.rs} states the 28 bits as fifteen row bits, ten column bits
and three bank bits. The wrapper takes the controller's address above
three lane bits, and the lane selects one word of the controller's
burst of eight.

On the board, \code{BoardRouter} sends the peripheral every address
that matches \code{0x4000\_0000} under the mask \code{0xc000\_0000}.
The cut to 28 bits removes the base, so the memory sees offsets from
\code{0x4000\_0000}.

\dstables{Ddr3Per}

\subsection*{Behaviour}

\begin{itemize}
\item The bridge takes one burst at a time and serves it a word at a
time at consecutive word addresses. A read's words go on the lines
one after another and its beats go back as they are acknowledged,
the last of them marked; a write goes a beat at a time as its beats
arrive and is answered once after the last.
\item The bridge offers a request with \code{cyc} and \code{stb}, and
the controller takes it in a cycle with no \code{stall}. The bridge
then holds \code{cyc} alone until \code{ack}, which brings the word
for a read. The answer goes on the link when the link has room.
\item Every Wishbone line the bridge drives is a register.
\item The wrapper repeats the word across the controller's burst of
eight with the mask clear on that one word alone, so a write touches
only its word; a read keeps its lane until the burst comes back and
picks its word out of it. One request is in flight at a time, and a
write is acknowledged the cycle after the controller takes it.
\item The controller makes the clocks. It takes the board's 200~MHz on
\code{sys\_clk}, which is also its delay reference, runs the memory
at 400~MHz, and hands out \code{ui\_clk}, the 100~MHz the design
runs on, with \code{ui\_rst} high until that clock is good. In the
netlist the controller's \code{i\_ui\_clk} is the design's clock,
which is that same clock fed back. \code{sys\_rst} is active high, a
pulse at power-on on the board.
\item \code{calib} is the controller's \code{o\_calib\_complete}.
\item In simulation the model stalls for \code{MODEL\_WARMUP}, 64
cycles, and then answers each request a cycle after it takes it.
\code{calib} rises when the warm-up has passed.
\end{itemize}

\subsection*{Verification}

\begin{itemize}
\item \code{//ddr3:ddr3\_test} runs \code{words\_go\_in\_and\_come\_back}.
Behind an \code{AxiHost} and an \code{AxiPer}, it writes words at
\code{0x4000\_0000}, \code{0x4000\_0104} and \code{0x47ff\_fffc}, the
first while the model still calibrates. It reads them back and checks
that \code{calib} did not rise before the warm-up.
\item \code{//ddr3:ddr3\_test} also runs
\code{the\_controller\_is\_instantiated\_and\_not\_written}. It checks
that the Verilog and the VHDL of \code{Ddr3Per} instantiate
\code{ddr3\_wb32}. The controller's clock must be on \code{clk}, the
board's clock and the clock it makes on the ports, and its pads on
the ports. The bridge's module must be
written and the controller's must not.
\item \code{//cpu/vreteno:board\_test} runs
\code{the\_memory\_test\_runs\_on\_the\_board}, a compiled program that
writes and reads the memory through the model and must print
\code{ddr3 ok}, and \code{the\_netlist\_holds\_the\_controller}.
\item \code{//ddr3/sim:wrapper\_test} simulates the wrapper and the
controller against the Micron model under Vivado's simulator:
calibrated, two words written at different lanes and read back. It
is manual. \code{//ddr3/sim:example\_test} is AMD's own proof of the
controller as configured, the example design's traffic generator
writing the memory and reading it back through it; manual too.
\item \code{//cpu/vreteno/board/sim:}\allowbreak\code{board\_test} simulates the whole
board, this controller included, against two Micron models. It is
manual.
\end{itemize}

\subsection*{Used by}

\code{Board}, as \code{ddr3} behind the tracker \code{pddr3}, at
\code{0x4000\_0000}.

\subsection*{Limits}

The Rust run of \code{Ddr3} is a model and not the controller. A
lowered netlist needs \code{//ddr3:hdl}, the wrapper, beside it, and
the controller's library, \code{//ddr3:ddr3\_mig}, which the build
generates under Vivado and which stays out of the tree, being AMD's.
The bridge serves one burst at a time. \code{//docs:vreteno} states
that the design with MIG as the controller has a bitstream that meets
timing, and that its run on the board is owed: the run it reports was
made with the controller before it.
